\documentclass[10pt,a4paper]{amsart}
\usepackage[margin=19mm]{geometry}
\usepackage{amsmath,amssymb,mathtools}
\usepackage[scr=rsfs,cal=euler]{mathalfa}
\usepackage{xfrac}
\usepackage[unicode,hidelinks,bookmarks=false]{hyperref}
\newcommand{\ie}{i.e.\ }
\newcommand{\cf}{cf.\ }
\newcommand{\resp}{respectively\ }
\newcommand{\Subsection}{\subsection}
\providecommand{\nobreakdash}{\nobreak\hbox{-}}
\setlength{\emergencystretch}{2em}
\multlinegap0pt
\usepackage{float}
\usepackage{amssymb}
\usepackage{tikz}
\newtheorem{lemma}{Lemma}
\newtheorem{proposition}{Proposition}
\usepackage{cleveref}
\crefname{lemma}{Lemma}{Lemmas}
\crefname{proposition}{Proposition}{Propositions}
\newcommand{\Ricci}{\mathrm{Ricci}}
\newcommand{\Uqu}{U_q(\mathfrak{u})}
\newcommand{\diffcalc}{\Omega_q}
\newcommand{\liftmap}{\ell}
\newcommand{\qalg}{\mathcal{B}}
\title{The Einstein condition for \\ quantum irreducible flag manifolds}
\author{Marco Matassa}
\date{Source: arXiv:2603.11786v1 (CC BY 4.0)}
\begin{document}
\maketitle

\noindent\emph{Source and licence.} The Einstein condition for \\ quantum irreducible flag manifolds. Version arXiv:2603.11786v1 (CC BY 4.0), distributed by arXiv under the Creative Commons Attribution 4.0 International licence, http://creativecommons.org/licenses/by/4.0/, as recorded in the arXiv licence metadata for this exact version and shown on the abstract page https://arxiv.org/abs/2603.11786v1. Full licence notice: \texttt{licenses/arxiv-2603.11786-CC-BY-4.0.txt}.\par\smallskip

\noindent\textit{Editorial scope.} Counted unit: the article's proposition prop:einstein-condition (source0.tex lines 757--772). The article's complete original proof of it is delivered with it (source0.tex lines 774--794, one \texttt{proof} environment of 21 lines). No mathematics is written, repaired or re-derived: the statement and the proof are the article's own lines, in the article's own notation, and no retained line is substituted or re-flowed.  Retained uncounted as the source-local context the proof needs: lines 496--500; lines 697--704.  Nothing was omitted from any retained span, so the package carries no omission record.  No retained line contains a citation, so the wrapper carries no bibliography and no bibliography entry.  No retained line contains a figure environment, an image-inclusion command or drawing code, so no image is delivered and the package declares no asset dependency.  Editorial formatting only, in addition: an AMS \texttt{amsart} preamble that restates the article's own declarations and macros used by the retained lines, namely the environments \texttt{lemma}, \texttt{proposition} on separate counters, together with the standard packages those lines need.  The retained environments print in this document as Lemma 1, Proposition 1 in order of appearance, whereas the article prints the same items with the numbering of its own full document.  The complete native source of the article as distributed in the arXiv e-print for this exact version is bundled unchanged as \texttt{sources/196237/source0.tex}.
\begin{lemma}
\label{lem:convex-lifting}
Suppose $\liftmap_1$ and $\liftmap_2$ are lifting maps.
Then $\liftmap = c_1 \liftmap_1 + c_2 \liftmap_2$ is a lifting map if and only if $c_1 + c_2 = 1$, in other words if $\liftmap$ is a convex combination of the two maps.
\end{lemma}

\begin{proposition}
\label{prop:ricci-proportional}
If $\liftmap$ is a $\Uqu$-equivariant lifting map then $\Ricci_\liftmap$ is a $\Uqu$-invariant element of $\diffcalc^1 \otimes_\qalg \diffcalc^1$. Furthermore, we have the following properties:
\begin{enumerate}
\item if $\liftmap_{+ -} : \diffcalc^{(1, 1)} \to \diffcalc^+ \otimes_\qalg \diffcalc^-$ then $\Ricci_{\liftmap_{+ -}} = a g_{+ -}$ for some coefficient $a$,
\item if $\liftmap_{- +} : \diffcalc^{(1, 1)} \to \diffcalc^- \otimes_\qalg \diffcalc^+$ then $\Ricci_{\liftmap_{- +}} = b g_{- +}$ for some coefficient $b$.
\end{enumerate}
\end{proposition}

\begin{proposition}
\label{prop:einstein-condition}
Let $\liftmap_{+ -}$ and $\liftmap_{- +}$ be two $\Uqu$-equivariant lifting maps as in \cref{prop:ricci-proportional}.
Denote by $a$ and $b$ the corresponding coefficients as defined there.

Suppose it holds that $a + b \neq 0$.
Then there exists a unique lifting map $\liftmap = c_1 \liftmap_{+ -} + c_2 \liftmap_{- +}$ such that the Einstein condition holds with non-zero Einstein constant, that is
\[
\Ricci_\liftmap = \lambda g, \quad \lambda \neq 0.
\]
In this case, the coefficients $c_1$ and $c_2$ for the lifting map $\liftmap$ are given by
\[
c_1 = \frac{b}{a + b}, \quad
c_2 = \frac{a}{a + b}.
\]
\end{proposition}

\begin{proof}
Note that, by \cref{lem:convex-lifting}, we have that $\liftmap = c_1 \liftmap_{+ -} + c_2 \liftmap_{- +}$ is a lifting map if and only if it is a convex combination of the two lifting maps $\liftmap_{+ -}$ and $\liftmap_{- +}$, which implies that $c_1 + c_2 = 1$.
It is clearly $\Uqu$-equivariant, since this is the case for both maps considered.

It immediately follows from the definition of the Ricci tensor that
\[
\Ricci_\ell = c_1 \Ricci_{\ell_{+ -}} + c_2 \Ricci_{\ell_{- +}}.
\]
Then, by \cref{prop:ricci-proportional}, there exist coefficients $a$ and $b$ such that
\[
\Ricci_\ell = c_1 a g_{+ -} + c_2 b g_{- +}.
\]
Hence this expression is proportional to the quantum metric $g = g_{+ -} + g_{- +}$ if and only if $c_1 a = c_2 b$.
Since $c_1 + c_2 = 1$, this is possible if and only if
\[
c_1 = \frac{b}{a + b}, \quad
c_2 = \frac{a}{a + b},
\]
which clearly requires that $a + b \neq 0$.
In this case, the Einstein condition holds with Einstein constant $\lambda = c_1 a = c_2 b$, which is non-zero by the assumption $a + b \neq 0$.
\end{proof}

\end{document}
